<<define "science">>
\subsection*{Citizen Radar}
\href{https://velocity.report}{velocity.report} is a citizen radar tool designed to help communities measure vehicle speeds with affordable, privacy-preserving Doppler sensors. It's built on a core physical truth: kinetic energy scales with the square of speed.
\par
\begin{minipage}{\linewidth}
\[ K_E = \tfrac{1}{2} m v^2 \]
\par
\begin{center}
where \(m\) is the mass and \(v\) is the velocity.
\end{center}
\end{minipage}
\par
\vspace{4pt}
A vehicle traveling at 40~mph has four times the crash energy of the same vehicle at 20~mph, posing exponentially greater risk to people outside the car. Even small increases in speed dramatically raise the likelihood of severe injury or death in a collision. By quantifying real-world vehicle speeds, \href{https://velocity.report}{velocity.report} produces evidence that exceeds industry standard metrics.
\par
\subsection*{Aggregation and Percentiles}
This system uses Doppler radar to measure vehicle speed by detecting frequency shifts in waves reflected from objects in motion. This shift (known as the \href{https://en.wikipedia.org/wiki/Doppler_effect}{Doppler effect}) is directly proportional to the object's relative velocity. When the sensor is stationary, the Doppler effect reports the true speed of an object moving toward or away from the radar.
\par
To structure this data, the \href{https://velocity.report}{velocity.report} application first records individual radar readings, then applies a greedy, local, univariate \emph{Time-Contiguous Speed Clustering} algorithm to group log lines into sessions based on time proximity and speed similarity. Each session, or ``transit,'' represents a short burst of movement consistent with a single passing object. This approach is efficient and reproducible, but in dense traffic or where objects overlap it may undercount vehicles by merging multiple objects into one transit.
\par
Undercounting can bias percentile metrics (like p85 and p98) downward, since fewer sessions can give disproportionate weight to slower vehicles. All reported statistics in this report are derived from these sessionised transits.
\par
Percentiles offer a structured way to interpret speed behaviour. The 85th percentile (p85) indicates the speed at or below which 85\% of vehicles traveled. The 98th percentile (p98) exceeds this industry-standard measure by capturing the fastest 2\% of vehicle speeds, providing a more robust view into trends among top speeders. By extending beyond p85, p98 identifies an additional 13\% of data that would otherwise be missed when trimming the top 15\%, offering clearer insight into high-risk driving patterns without letting single anomalous readings dominate.
\par
However, percentile metrics can be unstable in periods with low sample counts. To reflect this, our charts flag low-sample segments in orange and suppress percentile points when counts fall below reliability thresholds (fewer than 50 samples per roll-up period).
\par
<<end>>
